% ======================================================================
%  PARTIE 2 -- LIO-SAM                                       (~3 min 15)
%  (pas d'état de l'art complet : uniquement ce qu'il faut pour
%   comprendre LIO-SAM et justifier son choix)
% ======================================================================
\section{LIO-SAM : SLAM LiDAR-inertiel}

% ----------------------------------------------------------------------
% À DIRE (~45 s)
%  - SLAM = localisation et cartographie simultanées : pour se localiser
%    il faut une carte, et pour faire une carte il faut savoir où l'on
%    est -> on estime les deux en même temps, avec les mêmes capteurs.
%  - Deux formulations :
%      * Online SLAM (filtres : EKF/UKF, filtre particulaire) : on n'estime
%        que la pose courante -> léger, temps réel, mais les erreurs
%        s'accumulent (dérive) sans correction possible du passé ;
%      * Full SLAM : on estime toute la trajectoire -> optimisation
%        globale, permet la fermeture de boucle (loop closure), qui
%        corrige rétroactivement toute la trajectoire. Plus coûteux.
%  - Le Full SLAM se représente par un graphe de facteurs : nœuds = poses
%    successives du robot (et amers), facteurs = contraintes issues des
%    mesures (odométrie, LiDAR...). Sous hypothèse de bruits gaussiens,
%    maximiser la probabilité a posteriori revient à minimiser une somme
%    d'erreurs quadratiques (moindres carrés), résolue efficacement par
%    GTSAM ou g2o.
%  - Figure : poses en bleu clair, amers en bleu foncé.
%  - Transition : LIO-SAM repose justement sur un graphe de facteurs.
% ----------------------------------------------------------------------
\begin{frame}{Prérequis : SLAM et graphe de facteurs}
  \begin{columns}[c]
    \begin{column}{0.5\textwidth}
      \textbf{SLAM} : se \alert{localiser} $\leftrightarrow$ \alert{cartographier}
      \vspace{0.25cm}
      \begin{itemize}
        \item \textbf{Online SLAM} -- filtres\\
              {\small $p(x_t, m \mid Z_t, U_t)$ : léger, mais \alert{dérive}}
        \item \textbf{Full SLAM} -- optimisation\\
              {\small $p(X_T, m \mid Z_T, U_T)$ : toute la trajectoire\\ $\Rightarrow$ \alert{loop closure} possible}
      \end{itemize}
      \vspace{0.3cm}
      \textbf{Graphe de facteurs}
      \begin{itemize}
        \item nœuds = \alert{poses} du robot
        \item facteurs = \alert{contraintes} capteurs
        \item moindres carrés : GTSAM, g2o
      \end{itemize}
    \end{column}
    \begin{column}{0.48\textwidth}
      \centering
      \includegraphics[width=\linewidth]{facteur.png}
      \legende{Exemple de graphe de facteurs (Dellaert, 2005)}
    \end{column}
  \end{columns}
\end{frame}

% ----------------------------------------------------------------------
% À DIRE (~1 min 15)
%  - LIO-SAM : LiDAR Inertial Odometry via Smoothing And Mapping
%    (Shan et al., IROS 2020).
%  - Entrées : LiDAR 3D + IMU 9 axes (+ GPS optionnel).
%    Sorties : odométrie haute fréquence, poses clés (keyframes) et carte
%    de points globale optimisée.
%  - Couplage SERRÉ LiDAR-IMU dans un graphe de facteurs (GTSAM) : les
%    nœuds x_i sont les états du robot, reliés par 4 types de facteurs :
%      1. Pré-intégration IMU (orange) : intègre les mesures IMU brutes
%         entre deux keyframes -> prédiction du mouvement à haute
%         fréquence ;
%      2. Odométrie LiDAR (vert) : recalage (scan matching) du scan
%         courant sur une carte locale construite avec les dernières
%         keyframes -> corrige la prédiction IMU ;
%      3. GPS (jaune), quand il est disponible : contrainte de position
%         absolue -> limite la dérive à long terme ;
%      4. Loop closure (noir) : relie deux états NON consécutifs (lieu
%         déjà visité, trouvé par recherche de proximité puis recalage)
%         -> corrige rétroactivement toute la trajectoire entre les deux.
%  - Les 3 premiers ne contraignent que des états voisins dans le temps ;
%    seule la loop closure corrige globalement.
% ----------------------------------------------------------------------
\begin{frame}{LIO-SAM : principe}
  \centering
  {\small \textbf{L}iDAR \textbf{I}nertial \textbf{O}dometry via \textbf{S}moothing \textbf{A}nd \textbf{M}apping
   \quad {\color{textgray}(Shan \emph{et al.}, IROS 2020)}\par}
  \vspace{0.2cm}
  \begin{tikzpicture}[font=\scriptsize]
    \node[capteur=3.3cm] (in) at (0,0) {\textbf{Entrées} : LiDAR 3D, IMU 9 axes, GPS {\tiny(option)}};
    \node[boite=3.9cm, draw=accent, thick] (gr) at (4.9,0) {\textbf{Graphe de facteurs} (GTSAM)\\ couplage \alert{serré} LiDAR -- IMU};
    \node[capteur=3.6cm] (out) at (9.9,0) {\textbf{Sorties} : odométrie haute fréquence, keyframes, carte globale};
    \draw[fleche] (in) -- (gr);
    \draw[fleche] (gr) -- (out);
  \end{tikzpicture}\par
  \vspace{0.15cm}
  \includegraphics[width=0.84\textwidth]{liosam2.png}\par
  \vspace{0.1cm}
  \facteur{fIMU}{Pré-intégration IMU}{prédiction haute fréquence}\hfill
  \facteur{fLidar}{Odométrie LiDAR}{scan-matching, correction}\hfill
  \facteur{fGPS}{GPS}{position absolue}\hfill
  \facteur{fLoop}{Loop closure}{correction globale}
\end{frame}

% ----------------------------------------------------------------------
% À DIRE (~45 s)
%  - Côté logiciel, LIO-SAM = 4 nœuds ROS (de gauche à droite sur la
%    figure) qui s'échangent des messages :
%      1. imuPreintegration : reçoit l'IMU brute et publie une odométrie à
%         la fréquence de l'IMU (simple intégration), utilisée comme
%         prédiction par les autres nœuds ;
%      2. imageProjection : reçoit le nuage de points, l'IMU et cette
%         odométrie ; calcule une estimée initiale, organise le nuage et
%         corrige sa distorsion due au mouvement du robot pendant le scan
%         (deskewing) ; publie un message cloud_info ;
%      3. featureExtraction : extrait les points caractéristiques (arêtes
%         et plans, comme LOAM) et republie cloud_info ;
%      4. mapOptimization : recale le nuage sur la carte locale ->
%         odométrie LiDAR, puis optimise le graphe de facteurs (odométrie
%         LiDAR + GPS + loop closure) et publie l'odométrie optimisée.
%  - Cette odométrie optimisée repart vers imuPreintegration, qui
%    ré-estime le biais de l'IMU : cette boucle entre les nœuds = le
%    couplage serré de LIO-SAM.
%  - Pour mon framework : j'utilise la sortie odométrie de LIO (et plus
%    tard la trajectoire corrigée en cas de loop closure).
% ----------------------------------------------------------------------
\begin{frame}{LIO-SAM : architecture logicielle}
  \begin{columns}[c]
    \begin{column}{0.64\textwidth}
      \centering
      \includegraphics[width=\linewidth]{liosam.png}
      \legende{4 nœuds ROS et leurs messages (Shan \emph{et al.}, 2020)}
    \end{column}
    \begin{column}{0.34\textwidth}
      \small
      \begin{enumerate}
        \item \texttt{imuPreintegration}\\ {\footnotesize\color{textgray} odométrie IMU, \alert{biais IMU}}
        \item \texttt{imageProjection}\\ {\footnotesize\color{textgray} \alert{deskew} du nuage}
        \item \texttt{featureExtraction}\\ {\footnotesize\color{textgray} arêtes \& plans}
        \item \texttt{mapOptimization}\\ {\footnotesize\color{textgray} scan-matching + \alert{graphe}}
      \end{enumerate}
      \vspace{0.2cm}
      {\small $\circlearrowright$ odométrie LiDAR\\
      \phantom{$\circlearrowright$} $\rightarrow$ biais IMU\\[3pt]
      = couplage \alert{serré}\par}
    \end{column}
  \end{columns}
\end{frame}

% ----------------------------------------------------------------------
% À DIRE (~30 s)
%  - Pourquoi LIO-SAM : il exploite directement les capteurs des robots
%    (LiDAR 3D + IMU), il a une loop closure NATIVE (indispensable pour
%    corriger la dérive sur de longs trajets), et un portage ROS 2 mature.
%  - Surtout : le formalisme du graphe de facteurs est commun avec les
%    travaux du laboratoire (thèse de Clovis Adeler : diagnostic et
%    nouvelles contraintes GPS dans des architectures multi-robots) ->
%    base capable d'accueillir les futures briques du projet
%    (contraintes GPS multi-robots, planification, navigation).
%  - En pratique j'ai utilisé LIORF, une variante de LIO-SAM qui accepte
%    une IMU 6 axes (sans magnétomètre) : indispensable pour tester sur
%    KITTI, qui ne fournit pas de données magnétométriques.
%    Dans la suite je dis "LIO" pour LIO-SAM/LIORF.
%  - Si question : LIORF retire aussi le module d'extraction de
%    caractéristiques (cf. le titre du dépôt GitHub liorf).
%  - Si question sur les alternatives (état de l'art du rapport) :
%    FAST-LIO2 / Faster-LIO / DLIO / Point-LIO : très bons en odométrie
%    mais pas de loop closure native ; RTAB-Map : polyvalent mais plus
%    coûteux ; ORB-SLAM3 / VINS-Fusion : caméra seule, pas de LiDAR ;
%    SuMa++ : sémantique LiDAR mais académique, sans ROS 2 ni
%    multi-robots.
% ----------------------------------------------------------------------
\begin{frame}{Pourquoi LIO-SAM ?}
  \begin{columns}[c]
    \begin{column}{0.56\textwidth}
      \begin{itemize}
        \item[\ok] Capteurs des robots : \alert{LiDAR 3D + IMU}
        \item[\ok] \alert{Loop closure} native
        \item[\ok] Portage \alert{ROS 2} mature
        \item[\ok] \alert{Graphe de facteurs} : même formalisme\\ que les travaux du laboratoire
        \item[\ok] Base \alert{évolutive} pour ANR-SOS\\
              {\footnotesize\color{textgray} contraintes GPS multi-robots, navigation\dots}
      \end{itemize}
    \end{column}
    \begin{column}{0.38\textwidth}
      \begin{tcolorbox}[carte, title=En pratique : LIORF]
        Variante de LIO-SAM\\[4pt]
        \alert{IMU 6 axes}\\ {\footnotesize (sans magnétomètre)}\\[4pt]
        $\Rightarrow$ compatible \alert{KITTI}\\[6pt]
        {\footnotesize\color{textgray} dans la suite : « LIO »}
      \end{tcolorbox}
    \end{column}
  \end{columns}
\end{frame}
